\subsection{Simple Rings}

PROPOSITION 1. --- Let A be a nonzero ring. The following conditions are equivalent:

\begin{enumerate}
\def\labelenumi{(\roman{enumi})}
\item
  The A-module \(A_{s}\) is isotypical.
\item
  The ring A is left Artinian, and every two-sided ideal of A is equal to 0 or A.
\item
  The ring A is left Artinian, and there exists a left A-module S that is simple and faithful.
\end{enumerate}

If these conditions are satisfied, then the A-module \(A_{s}\) has finite length and is isotypical of type S, and every simple A-module is isomorphic to S.

Let us prove that (i) implies (ii). Assume that (i) is satisfied. Then the finitely generated A-module \(A_{s}\) is semisimple, hence has finite length and is Artinian (VIII, p.~71, Proposition 10); consequently, the ring A is left Artinian. The endomorphisms of the left A-module \(A_{s}\) are the right multiplications by the elements of A. Since the left A-module \(A_{s}\) is isotypical, it follows from Proposition 6, b) of VIII, p.~86 that the (A, A)-bimodule \(_{s}A_{d}\) is simple. The sub-bimodules of \(_{s}A_{d}\) are the two-sided ideals of A, so (i) implies (ii).

Let us prove that (ii) implies (iii). The ring A is not reduced to 0; consequently, there exists a simple A-module S. The annihilator of S is a proper two-sided ideal of A. If we assume that (ii) is satisfied, then that annihilator is equal to 0. The A-module S is then faithful, and (ii) implies (iii).

Let us prove that (iii) implies (i). Assume that (iii) is satisfied. Then there exists an integer \(m \geqslant 1\) such that \(A_{s}\) is isomorphic to a submodule of \(S^{m}\) (VIII, p.~50, Proposition 5, a)). Since \(S^{m}\) is an isotypical A-module of type S, the same holds for \(A_{s}\) (VIII, p.~61, Proposition 2); hence, (iii) implies (i).

Suppose that conditions (i) through (iii) are satisfied. We have seen above that the A-module \(A_{s}\) has finite length and is isotypical of type S. Every simple left A-module is isomorphic to a quotient of \(A_{s}\) , hence to S.

DEFINITION 1. --- We say that a ring A is simple if it satisfies the equivalent conditions (i), (ii), and (iii) of Proposition 1. Let K be a commutative field; a K-algebra is called simple if its underlying ring is simple.

Remarks. --- 1) Recall that by Theorem 1 of II, §9, No.~1, p.~115, the following properties are equivalent:

\begin{enumerate}
\def\labelenumi{(\roman{enumi})}
\item
  The A-module \(A_{s}\) is simple.
\item
  The ring A is not reduced to 0, and there are no left ideals of A distinct from 0 or A.
\item
  The ring A is a field. Consequently, by Proposition 1 (condition (ii)), commutative simple rings are nothing but commutative fields.
\end{enumerate}

\pandocbounded{\includegraphics[keepaspectratio]{images/part001_205e3f38.jpg}}

\begin{enumerate}
\def\labelenumi{\arabic{enumi})}
\setcounter{enumi}{1}
\tightlist
\item
  We sometimes say that a ring A is quasi-simple if it is not reduced to 0 and if its only two-sided ideals are 0 and A. We say that A is primitive if it admits a faithful simple module. By Proposition 1, every simple ring is quasi-simple. Since every nonzero ring admits a simple module and the annihilator of a simple module is a two-sided ideal, we see that every quasi-simple ring is primitive. However, there exist quasi-simple rings that are not simple and primitive rings that are not quasi-simple (VIII, p.~128, Exercise 2); such rings are not left Artinian.
\end{enumerate}

THEOREM 1 (Wedderburn). --- A ring is simple if and only if it is isomorphic to a matrix ring \(\mathbf{M}_r(\mathrm{D})\) , where \(r \geqslant 1\) is an integer and D a field.

Lemma 1. --- Let A be a simple ring, S a simple left A-module, and D the opposite ring of the field \(\operatorname{End}_{\mathrm{A}}(\mathrm{S})\) . Then S is an invertible (A,D)-bimodule. It is also a finite-dimensional right vector space over D, and the mapping \(a \mapsto a_{S}\) is a ring isomorphism from A to \(\operatorname{End}_{\mathrm{D}}(\mathrm{S})\) .

By Proposition 1, the A-module \(A_{s}\) has finite length and is isotypical of type S. Hence, there exists an integer \(m \geqslant 1\) such that the A-modules \(A_{s}\) and \(S^{m}\) are isomorphic. Then the A-module S is projective and finitely generated. It is generating (VIII, p.~80, Theorem 1), and Lemma 1 follows from Theorem 1 of VIII, p.~101, (ii)⇒(i) and (ii)⇒(iii) applied to the \((\mathrm{A},\mathrm{D})_{\mathbf{Z}}\) -bimodule S.

Lemma 2. --- Let D be a field and V a right vector space of finite dimension \(r \geqslant 1\) over D. Then V is a simple module over the ring \(\mathrm{E} = \mathrm{End}_{\mathrm{D}}(\mathrm{V})\) , and its commutant is equal to \(D_{V}\) . The ring E is simple, and its left length is equal to r.

We know that V is a simple E-module (VIII, p.~45, Example 3) and that its commutant is equal to \(D_{V}\) (VIII, p.~82, Corollary 1). Let \((x_{i})_{1\leqslant i\leqslant r}\) be a basis of V over the field D. The mapping \(u\mapsto(u(x_{i}))_{1\leqslant i\leqslant r}\) is an isomorphism from the E-module \(E_{s}\) to the E-module \(V^{r}\) ; consequently, the E-module \(E_{s}\) is isotypical of length r, so the ring E is simple.

Let us now prove Theorem 1. Recall (II, §10, No.~7, p.~349) that the ring \(\mathbf{M}_r(\mathrm{D})\) can be identified with the endomorphism ring of the right D-vector space \(\mathrm{D}_d^r\) ; moreover, every right vector space of finite dimension \(r\) over a field D is isomorphic to \(\mathrm{D}_d^r\) (II, §7, No.~1, p.~292). Theorem 1 therefore follows from Lemmas 1 and 2.

Remark 3. --- Let A be a simple ring, S a simple A-module, and D the opposite ring of the field \(\operatorname{End}_{\mathrm{A}}(\mathrm{S})\) . Then the A-module \(A_{s}\) has finite length, and \(\dim_{\mathrm{D}}(\mathrm{S})\) is equal to \(\operatorname{long}(\mathrm{A})\) . Indeed, by Lemma 1, the ring A is isomorphic to \(\operatorname{End}_{\mathrm{D}}(\mathrm{S})\) ; we then apply Lemma 2.

COROLLARY 1. --- a) The center of a simple ring is a field.

\begin{enumerate}
\def\labelenumi{\alph{enumi})}
\setcounter{enumi}{1}
\item
  The opposite ring of a simple ring is simple.
\item
  The left length of a simple ring is equal to its right length.
\end{enumerate}

Let D be a field, Z its center, and V a right vector space over D of finite dimension \(r \geqslant 1\) . We denote the endomorphism ring of V by E.

The mapping \(z \mapsto z_{V}\) is an isomorphism from Z to the center of E by Corollary 2 of VIII, p.~83. Assertion a) follows. The dual \(V^{*}\) of V is a right vector space over the opposite field \(D^{o}\) of D, and its dimension is equal to r. The mapping \(u \mapsto {}^{t}u\) is an isomorphism from the opposite ring \(E^{o}\) of E to the ring \(\operatorname{End}_{\mathrm{D}^{\circ}}(\mathrm{V}^{*})\) . Consequently, the ring \(E^{o}\) is simple, and the rings E and \(E^{o}\) have the same left length, equal to r (Lemma 2).

COROLLARY 2. --- Let r and \(r'\) be strictly positive integers, and let D and \(D'\) be fields. The rings \(\mathbf{M}_{r}(\mathrm{D})\) and \(\mathbf{M}_{r'}(\mathrm{D}')\) are isomorphic if and only if we have \(r = r'\) and the fields D and \(D'\) are isomorphic.

The condition is obviously sufficient.

Conversely, suppose that the rings \(\mathrm{B} = \mathbf{M}_{r}(\mathrm{D})\) and \(\mathrm{B}' = \mathbf{M}_{r'}(\mathrm{D}')\) are isomorphic. Since r is the length of \(B_{s}\) and \(r'\) that of \(B'_{s}\) (Lemma 2), we have \(r = r'\) . Moreover, B is Morita equivalent to D and \(B'\) to \(D'\) (VIII, p.~102,

Example 1). Consequently, the fields D and D' are Morita equivalent, hence isomorphic (VIII, p.~111, Proposition 13, c)).

COROLLARY 3. --- Let K be a commutative field, and let A be a K-algebra of finite degree with simple underlying ring. There exist an integer r and a K-algebra D of finite degree over K that is a field such that A is isomorphic to \(\mathrm{M}_{r}(\mathrm{D})\) . In particular, if K is algebraically closed, then A is isomorphic to a matrix algebra over K.

Let S be a simple left A-module; it is a finite-dimensional K-vector space over K. Its commutant is therefore an algebra of finite degree over K. The first assertion then follows from Lemma 1. If K is algebraically closed, then D = K by Theorem 1 of VIII, p.~47.

Remark 4. --- Let K be an algebraically closed commutative field, and let A be an algebra of finite degree over K. The algebra A is simple if and only if there exists an integer \(n \geqslant 1\) such that A is isomorphic to \(\mathrm{M}_{n}(\mathrm{K})\) . Its center is then isomorphic to K.

